\begin{claim}{D.2.20}{Proved}{The ninth-stage target tests and an ordinary extension}
Retain the notation of D.2.17--D.2.19, and let
$F=\fib(T\to K(\FF_2,2))$. The stable group and the
actual mod-two Hurewicz map satisfy
\[
 \pi_8\Sigma^\infty F\cong(\ZZ/2)^3,
 \qquad
 h:\pi_8F\longrightarrow H_8(F;\FF_2)
 \text{ is injective}.
 \tag{D.2.20h1}
\]
If $\overline\alpha:S^5\to T$ is the class of D.2.18
and $\lambda:S^7\to T$ is the generator of D.2.19, then
\[
 \overline\alpha\nu_5=0,
 \qquad \lambda\eta_7=0
 \quad\text{in }\pi_8T.
 \tag{D.2.20h2}
\]
The ordinary map of D.2.19 can be chosen to extend over
a finite nine-skeleton of a CW model of
$A=Q\Sigma K(\FF_2,2)$, preserving the homology values
of D.2.17. The statement concerns this finite stage;
it does not settle extension over the full source or
the ordinary nullification test in D.2.4.
\end{claim}
\begin{proof}
\textbf{Every degree-nine fiber class.}
Keep the fiber notation
\[
 A_4=a^2,\quad z_5=e+ab,\quad u_5=\xi a,
 \quad k=\{a,b\},\quad R=\{a,c\},
\]
and the degree-seven classes $P,Q$ of D.2.18.
The degree-eight fiber basis of D.2.19 is
\[
 W_8,C_8,L_8,M_8,J_1,J_2,J_3.
\]
Let $m,n,o,t\in H_9B$ be the labels dual to
$u^3v,u^2w,v^3,\operatorname{Sq}^4w$, respectively.
The full free double-loop basis of $H_9T$ has twenty-six
elements:
\[
\begin{gathered}
 m,n,o,t,\\
 ag,ah,a\xi b,aR,bf,bx,bk,cd,ce,c\xi a,\\
 a^2d,a^2e,a^2\xi a,abc,b^3,a^3b,\\
 \{a,f\},\{a,x\},\{b,d\},\{b,e\},\{a,k\},\xi c.
\end{gathered}
 \tag{D.2.20h3}
\]
The restricted Browder identity identifies
$\{b,\xi a\}$ with $\{a,k\}$; it is not an extra
generator. The base contributes dimension four and
the mixed terms in $H_9(B\times F)$ contribute
$2,2,3,4$. Hence $\dim H_9F=11$.
A complete fiber basis is
\[
\begin{split}
 Z_A&=A_4z_5,\qquad U_A=A_4u_5,\\
 N_9&=\xi c+aR+A_4u_5,\\
 K_1&=\{a,f\}+aR,\qquad
 K_2=\{a,x\}+bk,\\
 K_3&=\{b,d\},\qquad
 K_4=\{b,e\}+bk,\qquad K_5=\{a,k\},\\
 P_m&=m+ag+ce+a^2e+bf+a^3b,\\
 P_n&=n+ah+cd+a^2d,\qquad
 P_o=o+bx+b^3.
\end{split}
 \tag{D.2.20h4}
\]
The first two are products of primitive fiber classes;
the other nine are primitive. The top-operation diagonal
that needs particular care is
\[
 \overline\Delta(\xi c)
 =u_5\otimes A_4+A_4\otimes u_5
       +R\otimes a+a\otimes R.
 \tag{D.2.20h5}
\]
Apply top-operation additivity and Cartan to
$\Delta c=c\otimes1+1\otimes c+a\otimes a$.
The two product corrections in $N_9$ cancel exactly
the four terms in \textup{(D.2.20h5)}.
For example, the reduced diagonal of $\{a,x\}$ is
$k\otimes b+b\otimes k$, canceled by $bk$;
the other primitive assertions follow by the same
finite coproduct calculation. These are the exact
restricted-operation identities of
\cite[Chapter~III, Theorem~1.3(2)--(5), pp.~217--218]
{I_CohenLadaMay1976}.

Every displayed class has zero image under $q_*$.
For the polynomial and Browder primitives this follows
directly: the only degree-nine primitive in $H_*B$
is the label $t$, and none of those expressions has
a $t$ coefficient after applying $q_*$. The fourth-square
calculation below gives
$\operatorname{Sq}_*^4N_9=u_5$, whose image under $q_*$
is zero, whereas $\operatorname{Sq}_*^4t=d\ne0$.
Thus $N_9$ also has zero possible $t$ coefficient.
The primitive-kernel argument from D.2.18 puts these
classes in the fiber factor, and the two products
plainly belong there. The dimension count proves
completeness.

\medskip\noindent
\textbf{The complete operation table, including the top correction.}
With $\beta=\operatorname{Sq}_*^1$, the relevant operations are
\[
\begin{array}{c|c|c|c}
 &\beta&\operatorname{Sq}_*^2&\operatorname{Sq}_*^4\\ \hline
 Z_A&W_8&0&0\\
 U_A&0&0&0\\
 N_9&0&R&u_5\\
 K_1&0&0&0\\
 K_2&J_1+J_2&R&0\\
 K_3&J_1&0&0\\
 K_4&J_2&0&0\\
 K_5&0&0&0\\
 P_m&W_8&0&0\\
 P_n&0&P+Q&z_5\\
 P_o&L_8+M_8&P+Q&z_5.
\end{array}
 \tag{D.2.20h6}
\]
The $N_9$ row uses the top Nishida formula, not its
version below the top operation. Here $\xi c=Q^5c$,
$|c|=4$, and the nonlinear partition $(0,2)$ gives
\[
 \operatorname{Sq}_*^2(\xi c)
 =Q^3c+Q^4\beta c+\{c,\operatorname{Sq}_*^2c\}
 =\{c,a\}=R.
 \tag{D.2.20h7}
\]
The first term vanishes by instability and $\beta c=0$.
For the fourth square the only nonzero binomial term
is the term indexed by two:
\[
 \operatorname{Sq}_*^4(\xi c)
 =Q^3(\operatorname{Sq}_*^2c)=Q^3a=\xi a.
 \tag{D.2.20h8}
\]
Both possible nonlinear partitions have a zero factor.
The added terms $aR+A_4u_5$ have zero second and fourth
squares. These coefficients follow directly from the
primary formula in
\cite[Chapter~III, Theorem~1.3(3), pp.~217--218]
{I_CohenLadaMay1976}.

For the polynomial rows, Cartan on the Serre generators gives
\[
\begin{split}
 \beta m&=\beta n=0,\qquad \beta o=j+l,\\
 \operatorname{Sq}_*^2m&=0,\qquad
 \operatorname{Sq}_*^2n=\operatorname{Sq}_*^2o=g+h,\\
 \operatorname{Sq}_*^4m&=0,\qquad
 \operatorname{Sq}_*^4n=\operatorname{Sq}_*^4o=e.
\end{split}
\]
For example
$\operatorname{Sq}^2(u^2v)=v^3+u^2w$
and $\operatorname{Sq}^4(uv)=u^2w+v^3$.
Consequently
\[
\begin{split}
 \beta P_m&=af+af+a^4=W_8,\\
 \beta P_o&=j+l+ax+b(d+e)+ab^2=L_8+M_8,\\
 \operatorname{Sq}_*^2P_n&=\operatorname{Sq}_*^2P_o=P+Q,\\
 \operatorname{Sq}_*^4P_n&=\operatorname{Sq}_*^4P_o=e+ab=z_5.
\end{split}
\]
The bracket entries follow from its Cartan rule;
in particular $\operatorname{Sq}_*^2K_2=R$.
Thus the top operation is not the only source of
$R$ in this degree. The unit formulas use
\cite[\S2, nos.~8--9]{I_Serre1953}.

\medskip\noindent
\textbf{All second Bocksteins needed for integral normalization.}
The first-Bockstein images in degree eight span
$W_8,J_1,J_2,L_8+M_8$. Its cycle space is spanned
by those four classes and $C_8$. Thus the second
page in degree eight is one-dimensional on $[C_8]$.
There are two essential differentials:
\[
 \beta_2[P_n]=[C_8],\qquad
 \beta_2[P_m+Z_A]=[C_8].
 \tag{D.2.20h9}
\]
We check both, since omitting the second changes the
homotopy-group presentation below.
On the second page for $T$, $[a]=0$ because
$\beta b=a$, so $[P_n]=[n]+[c][d]$.
The unit $B\to T$ preserves the known differential
$\beta_2[d]=[c]$ belonging to $H_4B=\ZZ/4$.
Multiplicativity for the Pontryagin product gives
$\beta_2([c][d])=[c]^2=[C_8]$.
The term $\beta_2[n]$ lies in the base factor,
whereas $P_n$ and $C_8$ are fiber classes.
The space splitting $T\simeq B\times F$ splits
the Bockstein spectral sequence as a tensor product
over $\FF_2$. Its positive-degree base and fiber
factors have zero intersection. Therefore
$\beta_2[n]=0$ and the first equality follows.

For the second, compute
\[
 P_m+Z_A=m+ag+ce+bf,
 \qquad \beta(bg)=ag+bf.
\]
Hence $[P_m+Z_A]=[m]+[c][e]$.
The equality $\beta x=d+e$ gives $[e]=[d]$.
The same multiplicativity and base--fiber argument
forces $\beta_2[m]=0$ and proves the second equality
in \textup{(D.2.20h9)}.
These formulas and the universal coefficient theorem give
\[
 H_8F=\ZZ/4\oplus(\ZZ/2)^4.
 \tag{D.2.20h10}
\]

For completeness, the entire first-Bockstein cycle
space in degree nine has basis
\[
 U_A,N_9,K_1,K_5,K_2+K_3+K_4,P_n,P_m+Z_A.
 \tag{D.2.20h11}
\]
The second Bockstein vanishes on the first five.
For three of them this follows already from
\[
 U_A=\beta(z_5u_5),\quad
 K_1=\beta(\{b,f\}+bR),\quad
 K_5=\beta\{b,k\},
\]
followed, where necessary, by the actual projection
from $T$ to $F$. We next give integral representatives
for the other two. Thus \textup{(D.2.20h9)} accounts
for every nonzero second Bockstein relevant to the
nine-cell boundary matrix.

\medskip\noindent
\textbf{Two specified classes have genuine integral cycle lifts.}
Choose integral label chains $a_2,b_3,c_4,y_5,x_6$
with $db=2a$, $dc=dy=0$, and $dx=2y$, where the
reduction of $y$ is $d+e$. The coefficient two is
the already computed exponent of $H_5B$; even
chain changes give these compatible label choices.
For a degree-one Browder chain operation the differential is
\[
 d\lambda(s,t)=-\lambda(ds,t)-(-1)^{|s|}\lambda(s,dt).
\]
It follows that $\lambda(a,x)-\lambda(b,y)$ is an
integral cycle. Its reduction is exactly
$K_2+K_3+K_4$, since the two $bk$ terms cancel.

Since $c$ is an integral cycle of even degree,
the configuration chain $e_1(c\otimes c)$ is an
integral cycle. The terms $aR$ and $a^2\xi a$
are also products of integral cycles, so their
sum gives an integral lift of $N_9$ in $T$.
The configuration-chain conventions are those of
\cite[Chapter~III, proof of Lemma~3.10, p.~232]
{I_CohenLadaMay1976}.
Apply the actual space projection $T\simeq B\times F\to F$
to these chains. Their mod-two classes are unchanged,
because they lie in the fiber factor. Both can
therefore be chosen as zero-integral-boundary nine-cells
in a parity-preserving stable cellular model.

The class $P_m+Z_A+P_n$ has zero first and second
Bocksteins. Since \textup{(D.2.20h10)} has exponent
at most four, the same filtered integral normalization
makes it a zero-boundary nine-cell. Its operation
vector is $(P+Q,z_5)$ in the second and fourth columns.
This additional relation is retained in the following
complete component calculation.

\medskip\noindent
\textbf{All stable degree-eight components and their actual maps.}
Work two-locally, and take a finite stable approximation to
$\Sigma^\infty F$ through dimension ten before cancelling odd-unit
cell pairs. Retain its nine-skeleton, with lower skeleta $L_i$.
This preserves the required stable group and removes extraneous
nine-cells killed by unit ten-cell incidences. The lower model is
the one in D.2.18--D.2.19:
\[
 L_6=C\vee M(\ZZ/4,5)\vee S^6_{\rho_0},
 \qquad
 C=\cofib\bigl((2,\eta_4):S^4\vee S^5\to S^4\bigr).
 \tag{D.2.20t1}
\]
All groups of the $L_i$ in this calculation are stable groups.
The quotient $C\to S^5\vee S^6$ injects $\pi_8C$ into
$\ZZ/8\{\nu_5\}\oplus\ZZ/2\{\eta_6^2\}$.
Indeed the preceding term is $\pi_4^S=0$, and the following
map is $(x,y)\mapsto2x+4y$ in $\ZZ/8$.
Its kernel is cyclic of order four. Choose its generator $\omega$
so that
\[
 \pi_8C=\ZZ/4\{\omega\},\qquad
 q_*\omega=(2\nu_5,\eta_6^2),\qquad
 \delta\eta_6^2=2\omega.
 \tag{D.2.20t2}
\]
The last equality follows from $q_*\delta=(\eta_5,0)$ and
$\eta^3=4\nu$ two-locally. The stable stems and Moore
coextension formulas used here are the opened formulas in
\cite[\S2, (2.1)--(2.2), p.~149]{I_Oka1986}.

For the other summands, the Moore exact sequence and the
order-two lift $b_0$ of D.2.19 give
\[
\begin{split}
 \pi_8M(\ZZ/4,5)
   &=\ZZ/4\{n_0=\iota_5\nu_5\}
       \oplus\ZZ/2\{b_0\eta_7\},\\
 \pi_8S^6_{\rho_0}&=\ZZ/2\{c_0\eta_7\},
       \qquad c_0=\rho_0\eta_6.
\end{split}
 \tag{D.2.20t3}
\]
Here $\iota_5$ is the bottom Moore inclusion, whose
homology class is $u_5$.
The top pinch sends $b_0\eta_7$ to $\eta_6^2$;
it has order two because $b_0$ does.
Also $t\eta_6^2=i\eta^3=0$ in the mod-four Moore summand.

The seven-cell attachments
\[
 \gamma+t,\quad \gamma+2\rho_0+\delta,\quad\delta,\quad0
\]
have $\eta^2$-images $0,2\omega,2\omega,0$.
Their maps on the preceding $\eta$-groups have kernel
generated by $\eta$ on the cell combinations $P+Q$ and $R$.
Thus the cofiber sequence gives
\[
\begin{split}
 0\longrightarrow&
 \ZZ/2\{\overline\omega\}\oplus\ZZ/4\{n_0\}
 \oplus\ZZ/2\{b_0\eta_7,c_0\eta_7\}\\
 \longrightarrow&\ \pi_8L_7
 \longrightarrow (\ZZ/2)^2\longrightarrow0.
\end{split}
 \tag{D.2.20t4}
\]
The $R$ class has the actual order-two lift
$\ell_0=\lambda\eta_7$, since that seven-cell is freely wedged.
For the other lift, change the seven-cell basis from $P,Q$ to
$P+Q,Q$. The new $P+Q$ attachment is $\gamma+2\rho_0$.
Choose a lift $\vartheta$ of its top $\eta$ in the subcomplex
\[
 D=(C\vee S^6_{\rho_0})
          \mathop{\cup}_{\gamma+2\rho_0}e^7_{P+Q}.
 \tag{D.2.20t5}
\]
It exists because $\gamma\eta_6=0$ and $2\rho_0\eta_6=0$.
The element $2\vartheta$ comes from $\pi_8(C\vee S^6)$;
in $\pi_8L_7$ it therefore belongs to the subgroup generated
by $\overline\omega$ and $c_0\eta_7$. In particular it has
no $n_0$ component. This specified representative excludes
an ambiguity that an abstract group extension alone would leave.

\medskip\noindent
\textbf{The precise $\kappa\eta$ relation.}
Set $C'=C\cup_\delta e^7_Q$. Equation
\textup{(D.2.20t2)} and $\delta\eta_6\ne0$ give
\[
 \pi_8C'=\ZZ/2\{\overline\omega\}.
\]
Choose $\kappa'\in\pi_7C'$ with top-cell degree two on $Q$.
In its cone the new degree-eight class has Bockstein $Q$.
Since $\operatorname{Sq}_*^2Q=z_5$ and
$\beta z_5=A_4$, the Adem identities force its
second-square vector to be $e_6$, the only six-class in $C'$.
The quotient of $C'$ by its bottom sphere is
\[
 C'/S^4\simeq C\eta_{5}\vee S^6_{e_6},
\]
where $C\eta_5$ has cells in dimensions five and seven.
The $S^6$ component of $\kappa'$ is therefore $\eta_6$.
Consequently $\kappa'\eta_7$ has a nonzero $\eta_6^2$
component in this quotient, and hence equals
$\overline\omega$.

The normalized $\kappa$ of D.2.19 has second-square vector
$d_6$. Both $\kappa$ and $\kappa'$ project to twice the
$Q$ cell. Their difference in $\pi_7L_7$ is therefore
torsion. The complete detectors of D.2.19 show that this
difference is $b_0$ plus possibly $v_0$: the required
second-square difference is $d_6+e_6$, and no $c_0$ term
is possible. As $v_0\eta_7=i\nu\eta=0$, we obtain
\[
 \kappa\eta_7=\overline\omega+b_0\eta_7.
 \tag{D.2.20t6}
\]
This calculation uses the fixed quotient normalization of
$\gamma$ and the fixed second-square normalization of
$\kappa$ from the preceding claims.

\medskip\noindent
\textbf{The entire group before the nine-cells.}
The eight-cell attachments of D.2.19, ordered as
$W_8,C_8,L_8,M_8,J_1,J_2,J_3$, are
\[
 0,\quad v_0,\quad\kappa,\quad
 \kappa+b_0+c_0,\quad b_0,\quad0,\quad2\lambda+b_0.
\]
Their $\eta$-images generate exactly
$\langle\kappa\eta_7,b_0\eta_7,c_0\eta_7\rangle$.
By \textup{(D.2.20t6)}, these kill
$\overline\omega,b_0\eta_7,c_0\eta_7$ in
\textup{(D.2.20t4)}. They also kill $2\vartheta$, by
the representative in \textup{(D.2.20t5)}, and impose
no relation on $n_0$. The resulting torsion subgroup is exactly
\[
 T_8=\ZZ/4\{n_0\}\oplus
        \ZZ/2\{\vartheta,\ell_0\}.
 \tag{D.2.20t7}
\]
The kernel on the following free degree-seven cell groups
is freely generated by
\[
 W_8,\quad J_2,\quad2C_8,\quad2J_1,\quad2(M_8-L_8).
\]
The extension by this free group splits. More invariantly,
write $D_8=M_8-L_8$ in integral cellular coordinates.
The integral Hurewicz map of $L_8$ has kernel precisely
$T_8$ and image
\[
 \Lambda=\ZZ_{(2)}\{W_8,J_2\}
       \oplus2\ZZ_{(2)}\{C_8,J_1,D_8\}
 \subseteq H_8(L_8;\ZZ_{(2)}).
 \tag{D.2.20t8}
\]
In particular the free lift of $2C_8$ will be denoted $a_0$.

For an attaching map with zero integral boundary, the complete
torsion detectors needed here are
\[
\begin{array}{c|c|c}
 &\operatorname{Sq}_*^2&\operatorname{Sq}_*^4\\ \hline
 n_0&0&u_5\\
 \vartheta&P+Q&z_5\\
 \ell_0&R&0.
\end{array}
 \tag{D.2.20t9}
\]
The $n_0$ row is the Hopf-invariant-one cone on the
Moore bottom, and the $\ell_0$ row is the cone of $\eta$
on the $R$ seven-sphere. For $\vartheta$, its top quotient
is $\eta$ on $P+Q$, giving its second square. Its chosen
representative in \textup{(D.2.20t5)} has no $u_5$ class,
so its fourth-square $u_5$ coefficient is zero. Finally,
dualize the Adem identity
$\operatorname{Sq}^2\operatorname{Sq}^3
=\operatorname{Sq}^5+\operatorname{Sq}^4\operatorname{Sq}^1$.
The top Bockstein is zero, while
$\operatorname{Sq}_*^3(P+Q)=A_4$ and $\beta z_5=A_4$.
Its fourth-square $z_5$ coefficient must therefore be one.
For the order-four class $n_0$, these primary detectors
see only the parity of its coefficient; we retain that
limitation in the next step.

\medskip\noindent
\textbf{Three zero-boundary relations kill the entire torsion subgroup.}
The genuine integral cycles $N_9$ and $K_2+K_3+K_4$ constructed
above can be chosen as zero-boundary nine-cells. By their
operation vectors and \textup{(D.2.20t9)}, their attachments
have the respective forms
\[
 \ell_0+(2r+1)n_0,\qquad \ell_0+2s n_0.
 \tag{D.2.20t10}
\]
Neither has a $\vartheta$ component. Subtracting the two
relations gives an odd multiple of $n_0$, so they kill
the entire order-four class $n_0$, and then $\ell_0$.
No value of an undetected even coefficient has been assumed.

The additional cycle $Z_0=P_m+Z_A+P_n$ has zero first and
second Bocksteins by \textup{(D.2.20h9)}. The integral
normalization explained above makes it a zero-boundary
nine-cell. Its operation vector is $(P+Q,z_5)$, so its
attachment is $\vartheta$ plus an even multiple of $n_0$,
with no $\ell_0$ term. It kills $\vartheta$ after the
preceding two relations. Thus all of $T_8$ is killed.
Both second Bocksteins in \textup{(D.2.20h9)} are essential
here; omitting the one on $P_m+Z_A$ would leave an incorrect
order-four stable summand.

\medskip\noindent
\textbf{The remaining quotient is the integral Hurewicz image.}
Let $\mathcal B_9$ be the image of the integral ninth
cellular boundary in $H_8L_8$. Every such boundary belongs
to the lattice $\Lambda$. Since the whole Hurewicz kernel
$T_8$ has already been killed by zero-boundary cells,
the complete stable homotopy quotient is exactly
\[
 \pi_8\Sigma^\infty F=\Lambda/\mathcal B_9.
 \tag{D.2.20t11}
\]
Two attachments with the same integral boundary differ
by $T_8$, so no free-lift extension remains in this formula.
In the integral homology calculation, the four
first-Bockstein images supply order-two lifts of
$W_8,J_1,J_2,D_8$, and $C_8$ is the order-four generator.
Consequently \textup{(D.2.20t8)} identifies the quotient
with the following subgroup of $H_8F$:
\[
 \pi_8\Sigma^\infty F
 \xrightarrow[\cong]{\ h_{\ZZ}\ }
 \langle2C_8,W_8,J_2\rangle
       \subseteq\ZZ/4\{C_8\}\oplus
           \ZZ/2\{W_8,J_1,J_2,D_8\}.
 \tag{D.2.20t12}
\]
Changing an order-two lift by $2C_8$ does not change
this subgroup. It is $(\ZZ/2)^3$. In particular, if
$a_0$ denotes the class with integral Hurewicz image
$2C_8$, then
\[
 \ker\bigl(h_{\FF_2}:\pi_8\Sigma^\infty F
                     \to H_8(F;\FF_2)\bigr)
       =\ZZ/2\{a_0\}.
 \tag{D.2.20t13}
\]
All nine-cell attachments have been accounted for:
the zero-boundary ones killed $T_8$, and their integral
boundary subgroup gives exactly \textup{(D.2.20t11)}.
Higher cells cannot change this stable group.

\medskip\noindent
\textbf{The ordinary Hurewicz map is injective.}
Write $F=\Omega^2 Z$. The double suspension--loop unit
\[
 F\longrightarrow\Omega^2\Sigma^2F
\]
has a homotopy retraction by the adjunction triangle
identity. Since $\Sigma^2F$ is five-connected, its
tenth homotopy group is stable. Therefore stabilization
gives a split injection
\[
 \pi_8F\lhook\joinrel\longrightarrow\pi_8\Sigma^\infty F.
 \tag{D.2.20t14}
\]
In particular the stable zeros $n_0=0$ and $\ell_0=0$
already prove the actual zeros in \textup{(D.2.20h2)}.

The remaining stable mod-two-Hurewicz-kernel element
$a_0$ is not an actual sphere map. Its integral
Hurewicz image is $2C_8$, with $C_8$ of order four.
In its stable cone a new degree-nine mod-two class $x_9$
satisfies
\[
 \beta x_9=C_8\pmod{\operatorname{im}\beta|_{H_9F}}.
\]
The old first-Bockstein images are
$W_8,J_1,J_2,L_8+M_8$; all have zero fourth square.
Since $\operatorname{Sq}_*^4C_8=A_4$, the Adem identity
$\operatorname{Sq}^1\operatorname{Sq}^4
=\operatorname{Sq}^5$ gives
\[
 \operatorname{Sq}_*^5x_9
   =\operatorname{Sq}_*^4\beta x_9=A_4\ne0.
 \tag{D.2.20t15}
\]
If $a_0$ came from an actual map $S^8\to F$, its
ordinary space cone would have this operation.
But $\operatorname{Sq}^5$ on degree-four cohomology
is zero by instability. This contradiction excludes
$a_0$ from the image of \textup{(D.2.20t14)}.
By \textup{(D.2.20t13)}, that image meets the stable
mod-two Hurewicz kernel only in zero. Hence
$\pi_8F\to H_8(F;\FF_2)$ is injective, as claimed.
No choice of a stable quadratic complement and no
identification of stable and ordinary eighth homotopy
groups has been used.

\medskip\noindent
\textbf{A source decomposition that treats every nine-cell at once.}
We now use the target conclusions to extend the ordinary map,
keeping the choices made through dimension eight explicit.
Put $Q_i=K_i/K_4$ as in D.2.18--D.2.19, and work with the
two-local quotient cell models there. The five eight-cell
attachments in \textup{(D.2.19w)} give, by cell slides,
\[
 Q_8\simeq D'\vee M(\ZZ/2,7)_F\vee M(\ZZ/4,7)_J
                  \vee C_1^8\vee S^8_s,
 \tag{D.2.20s1}
\]
where
\[
\begin{split}
 D&=\operatorname{cofib}\left(
  \begin{pmatrix}4&\eta_5\\0&2\end{pmatrix}:
  S^5\vee S^6\longrightarrow S^5\vee S^6\right),\\
 D'&=D\cup_{\tau_d+\tau_e}e^8_H,\qquad
 C_1^8=\operatorname{cofib}\bigl((2,\eta_6):S^6\vee S^7\to S^6\bigr).
\end{split}
 \tag{D.2.20s2}
\]
Indeed slide the cell $G_8$ over $H_8$, leaving the
degree-two attachment on $F$; slide $B_0^2$ over $A_0C$,
leaving the degree-four attachment on $J$; and slide
$Q^5A_0$ over $A_0C$, leaving the free eight-sphere.
The integral orientation of the last sphere is
$s=Q^5A_0-A_0C$. The cell $A_0C$ supplies the top
cell of $C_1^8$, while $H_8$ supplies the top cell of $D'$.
These are actual comparisons for the mapping problem:
the quotients are four-connected, their dimensions are
at most eight, and Freudenthal realizes the stable cell
slides unstably. Maps into the two-local space $T$ may
then be pulled back to the original finite skeleta.

In D.2.19 the correction is zero on the source six-sphere
$z$ and on the free seven-spheres $F,J$. Its nullhomotopies
on the eight-cells may therefore be chosen so that the
correction is constant on all the factors of
\textup{(D.2.20s1)} except $D'$. This makes its restrictions
to the two Moore factors, $C_1^8$, and the last sphere
literally zero. On $D'$ it retains the Moore extension
$\mu$, the bottom class $\overline\alpha$, and the
six-sphere value $\rho$. Its extension across the
$H$-cell can be any of the allowed nullhomotopies;
we will not assume that its effect on the next
homotopy group is an arbitrary free choice.

\medskip\noindent
\textbf{Every class in the remaining attaching group has zero
mod-two Hurewicz image.}
In $D$, the two classes $\tau_d,\tau_e$ are independent
order-two elements, by the complete calculation
\textup{(D.2.19p)}. Thus $\tau_d+\tau_e$ is nonzero
and has order two. The cofiber sequence
\[
 S^7\xrightarrow{\tau_d+\tau_e}D\longrightarrow D'
             \xrightarrow{p_H}S^8
 \tag{D.2.20s3}
\]
gives
\[
 \operatorname{im}\bigl((p_H)_*:\pi_8D'\to\pi_8S^8\bigr)
                  =2\ZZ_{(2)}.
 \tag{D.2.20s4}
\]
This is first computed in the stable exact sequence;
four-connectivity identifies all the relevant groups
with their ordinary counterparts. The complex $D$
has dimension seven and $H_7(D;\ZZ_{(2)})=0$:
its top cellular differential has nonzero degree two
on the free six-cell. The homology sequence of
\textup{(D.2.20s3)} consequently identifies
$H_8(D';\ZZ_{(2)})$ with $\ZZ_{(2)}$, by the
pinch onto the $H$-cell. It follows from
\textup{(D.2.20s4)} that every actual sphere map
$S^8\to D'$ has even degree on this sole integral
eight-dimensional homology generator. Therefore
\[
 h_{\FF_2}:\pi_8D'\longrightarrow H_8(D';\FF_2)
                  \quad\text{is zero}.
 \tag{D.2.20s5}
\]
There is no additional Tor summand in this last
homology group, since $H_7D'=0$. This argument
includes the entire torsion subgroup of $\pi_8D'$;
it does not infer its vanishing from incomplete
primary-square detectors.

\medskip\noindent
\textbf{The correction kills all actual nine-cell attachments.}
Let $f:D'\to T$ be the correction, with any of its
allowed fillings of the $H$-cell. It lifts to
$F=\fib(T\to B)$, since
$\Map_*(D',K(\FF_2,2))$ is weakly contractible.
Naturality and \textup{(D.2.20s5)} make the image
of every element of $\pi_8D'$ have zero mod-two
Hurewicz in $F$. The target result proved above
makes the actual Hurewicz homomorphism on $\pi_8F$
injective. Thus
\[
 f_*:\pi_8D'\longrightarrow\pi_8T
                       \quad\text{is zero}.
 \tag{D.2.20s6}
\]
The correction is constant on the other wedge
factors in \textup{(D.2.20s1)}. There are no cross
Whitehead terms in dimension eight by their
connectivities. Consequently its homomorphism
on all of $\pi_8Q_8$ is zero. Every nine-cell
attaching map has zero image, regardless of any
component not measured by primary operations.
Choose nullhomotopies on these cells to extend
the correction to $Q_9$.

Finally, the canonical summand already constructed
in \textup{(D.2.19y)--(D.2.19z)} extends through
source relative dimension nine. We keep that same
map and add the pullback of the correction along
$K_9\to Q_9$. This gives an actual ordinary map
$G_9:K_9\to T$ extending the allowed $G_8$ and
retaining every six-dimensional value in D.2.17.
The canonical summand has only been supplied
through dimension nine, and no extension over
the tenth or later cells follows from this proof.

\end{proof}
